
\documentclass[aspectratio=169,10pt]{beamer}
\usepackage[T1]{fontenc}
\usepackage[utf8]{inputenc}
\usepackage[indonesian]{babel}
\usepackage{amsmath,amssymb}
\usepackage{tikz}
\usetikzlibrary{arrows.meta,positioning,calc}
\usepackage{booktabs}
\usepackage{array}
\usepackage{listings}
\usepackage{xcolor}
\usepackage{ragged2e}

\definecolor{pidknavy}{HTML}{123C63}
\definecolor{pidkblue}{HTML}{E8F3FB}
\definecolor{pidkteal}{HTML}{3BA7A9}
\definecolor{pidkgray}{HTML}{263238}
\definecolor{codebg}{HTML}{F7F9FB}
\definecolor{codeframe}{HTML}{BAC8D3}
\definecolor{coderule}{HTML}{0B5D89}

\usetheme{default}
\setbeamertemplate{navigation symbols}{}
\setbeamercolor{frametitle}{fg=white,bg=pidknavy}
\setbeamertemplate{frametitle}{\nointerlineskip\begin{beamercolorbox}[wd=\paperwidth,ht=0.78cm,dp=0.24cm,leftskip=0.28cm]{frametitle}\usebeamerfont{frametitle}\insertframetitle\end{beamercolorbox}}
\setbeamerfont{frametitle}{size=\Large,series=\bfseries}
\setbeamerfont{normal text}{size=\normalsize}
\setbeamertemplate{footline}{%
  \ifnum\value{framenumber}=1\relax\else%
  \leavevmode\hbox{%
  \begin{beamercolorbox}[wd=.36\paperwidth,ht=2.7ex,dp=1ex,leftskip=0.25cm]{footA}Pemrograman untuk Ilmu Data dan Komputasi\end{beamercolorbox}%
  \begin{beamercolorbox}[wd=.48\paperwidth,ht=2.7ex,dp=1ex,center]{footB}Program Studi Matematika ITERA\end{beamercolorbox}%
  \begin{beamercolorbox}[wd=.16\paperwidth,ht=2.7ex,dp=1ex,rightskip=0.25cm plus1fil]{footC}\hfill\insertframenumber/\inserttotalframenumber\end{beamercolorbox}}%
  \fi}
\setbeamercolor{footA}{bg=pidknavy,fg=white}
\setbeamercolor{footB}{bg=pidknavy,fg=white}
\setbeamercolor{footC}{bg=pidknavy,fg=white}

\lstdefinestyle{pidk}{
  language=Python,
  basicstyle=\ttfamily\footnotesize,
  keywordstyle=\color{coderule}\bfseries,
  commentstyle=\color{pidkteal},
  stringstyle=\color{teal!60!black},
  showstringspaces=false,
  backgroundcolor=\color{codebg},
  frame=single,
  rulecolor=\color{codeframe},
  framerule=0.45pt,
  columns=fullflexible,
  keepspaces=true,
  breaklines=true,
  aboveskip=3pt,
  belowskip=3pt
}
\lstset{style=pidk}

\newenvironment{pidkbox}[1]{\begin{block}{#1}\justifying}{\end{block}}
\setbeamercolor{block title}{bg=pidknavy,fg=white}
\setbeamercolor{block body}{bg=pidkblue,fg=pidkgray}
\setbeamerfont{block title}{series=\bfseries}

\title{Pertemuan 2\\Fungsi, Subprogram, dan Rekursi}
\subtitle{Pemrograman untuk Ilmu Data dan Komputasi\\MA25-21008 Kelas RA dan RB}
\author{Triyana Muliawati, S.Si., M.Si.\\Rifky Fauzi}
\institute{Program Studi Matematika\\Institut Teknologi Sumatera}
\date{Semester Ganjil 2026/2027}

\begin{document}

% PATCH M02 FINAL MARKER: TRIYANA-FIRST SUBPROGRAM RECURSION NO-RUJUKAN NO-BATAS

\begin{frame}[plain]
\begin{columns}[T,onlytextwidth]
\begin{column}{.47\linewidth}
\vspace{0.9cm}
\begin{tikzpicture}[scale=1]
\fill[pidknavy] (0,0) rectangle (5.7,7.0);
\foreach \i in {1,...,45}{\fill[cyan!55!white,opacity=.35] ({rnd*5.6},{rnd*6.8}) circle (0.018+rnd*0.035);}
\draw[cyan!45!white,opacity=.65,line width=.4pt] plot[smooth] coordinates{(.15,1.1)(1.1,1.5)(2.1,.8)(3.0,1.8)(4.2,1.2)(5.4,2.0)};
\draw[cyan!35!white,opacity=.45] (0.4,5.8)--(1.8,6.4)--(2.7,5.5)--(3.7,6.2)--(5.1,5.4);
\foreach \p in {(0.4,5.8),(1.8,6.4),(2.7,5.5),(3.7,6.2),(5.1,5.4)}{\fill[cyan!70!white] \p circle (.055);}
\node[align=center,text=white!72,font=\Large] at (3.0,3.8){$y=f(x)$};
\node[align=center,text=white!55,font=\normalsize] at (3.0,3.1){input $\to$ proses $\to$ output};
\end{tikzpicture}
\end{column}
\begin{column}{.50\linewidth}
\vspace{0.82cm}
{\Huge\bfseries\color{pidknavy} Pertemuan 2}\par\vspace{1mm}
{\LARGE\bfseries\color{pidknavy} Fungsi, Subprogram, dan Rekursi}\par\vspace{7mm}
{\large Pemrograman untuk Ilmu Data dan Komputasi}\par\vspace{2mm}
{MA25-21008 \quad Kelas RA dan RB}\par\vspace{6mm}
{Triyana Muliawati, S.Si., M.Si.}\par
{Rifky Fauzi}\par\vspace{6mm}
{Program Studi Matematika}\par
{Institut Teknologi Sumatera}\par\vspace{4mm}
{Semester Ganjil 2026/2027}
\vfill
\begin{pidkbox}{Hari ini}
Fungsi sebagai subprogram, input-output, scope, komposisi fungsi, dan rekursi.
\end{pidkbox}
\end{column}
\end{columns}
\end{frame}

\begin{frame}{Dari M1 ke M2}
\begin{columns}[T,onlytextwidth]
\begin{column}{.48\linewidth}
\begin{pidkbox}{Yang sudah dipakai di M1}
\begin{itemize}
\item ekspresi, tipe data, dan variabel;
\item input/output;
\item percabangan \texttt{if};
\item pengulangan \texttt{for} dan \texttt{while}.
\end{itemize}
\end{pidkbox}
\end{column}
\begin{column}{.48\linewidth}
\begin{pidkbox}{Yang mulai dibuat rapi di M2}
\begin{itemize}
\item potongan logika diberi nama;
\item input dan output dibuat eksplisit;
\item program dapat dipanggil ulang;
\item kesalahan lebih mudah dilacak.
\end{itemize}
\end{pidkbox}
\end{column}
\end{columns}
\vspace{3mm}
\begin{center}
\large Program yang baik bukan hanya memberi jawaban, tetapi juga menjelaskan alur berpikirnya.
\end{center}
\end{frame}

\begin{frame}{Apa itu subprogram?}
\begin{pidkbox}{Makna umum}
Subprogram adalah bagian program yang diberi nama untuk menjalankan tugas tertentu. Dalam Python, bentuk subprogram yang paling sering dipakai adalah \textbf{fungsi}.
\end{pidkbox}
\vspace{2mm}
\begin{columns}[T,onlytextwidth]
\begin{column}{.48\linewidth}
\begin{pidkbox}{Tanpa subprogram}
\begin{itemize}
\item satu program panjang;
\item banyak bagian berulang;
\item sulit diuji bagian demi bagian;
\item sulit dipakai ulang pada data lain.
\end{itemize}
\end{pidkbox}
\end{column}
\begin{column}{.48\linewidth}
\begin{pidkbox}{Dengan subprogram}
\begin{itemize}
\item satu tugas diberi nama;
\item input dan output tampak jelas;
\item dapat diuji terpisah;
\item dapat disusun menjadi program besar.
\end{itemize}
\end{pidkbox}
\end{column}
\end{columns}
\end{frame}

\begin{frame}{Fungsi sebagai mesin kecil}
\begin{center}
\begin{tikzpicture}[node distance=1.25cm,>=Stealth]
\node[draw,rounded corners,fill=pidkblue,minimum width=2.2cm,minimum height=1.0cm] (input) {Input};
\node[draw,rounded corners,fill=pidknavy!88,text=white,minimum width=3.2cm,minimum height=1.2cm,right=of input] (fungsi) {Fungsi};
\node[draw,rounded corners,fill=pidkblue,minimum width=2.2cm,minimum height=1.0cm,right=of fungsi] (output) {Output};
\draw[->,line width=1.2pt,color=pidknavy] (input)--(fungsi);
\draw[->,line width=1.2pt,color=pidknavy] (fungsi)--(output);
\node[below=.5cm of input,align=center] {$x$ atau data};
\node[below=.5cm of fungsi,align=center] {aturan / algoritma};
\node[below=.5cm of output,align=center] {$f(x)$ atau hasil};
\end{tikzpicture}
\end{center}
\begin{pidkbox}{Cara membacanya}
Fungsi menerima nilai, melakukan proses yang jelas, lalu mengembalikan hasil. Untuk mahasiswa Matematika, fungsi Python dapat dipandang sebagai bentuk komputasional dari transformasi matematika.
\end{pidkbox}
\end{frame}

\begin{frame}[fragile]{Anatomi fungsi Python}
\begin{lstlisting}
def nama_fungsi(parameter_1, parameter_2):
    langkah_1
    langkah_2
    return hasil
\end{lstlisting}
\begin{columns}[T,onlytextwidth]
\begin{column}{.48\linewidth}
\begin{pidkbox}{Saat mendefinisikan}
\begin{itemize}
\item \texttt{def}: mulai mendefinisikan fungsi;
\item nama fungsi: sebaiknya jelas;
\item parameter: variabel lokal di dalam fungsi.
\end{itemize}
\end{pidkbox}
\end{column}
\begin{column}{.48\linewidth}
\begin{pidkbox}{Saat memanggil}
\begin{itemize}
\item argumen: nilai nyata yang dikirim;
\item \texttt{return}: hasil dikirim kembali;
\item setelah \texttt{return}, fungsi selesai.
\end{itemize}
\end{pidkbox}
\end{column}
\end{columns}
\end{frame}

\begin{frame}[fragile]{Contoh 1: fungsi dari rumus matematika}
\begin{columns}[T,onlytextwidth]
\begin{column}{.42\linewidth}
\[
 p(x)=2-3x+x^2
\]
\[
 p(4)=2-12+16=6
\]
\end{column}
\begin{column}{.55\linewidth}
\begin{lstlisting}
def polinom_sederhana(x):
    return 2 - 3*x + x**2

hasil = polinom_sederhana(4)
print(hasil)
\end{lstlisting}
\end{column}
\end{columns}
\begin{pidkbox}{Perhatikan}
Nilai \texttt{x} pada definisi fungsi adalah parameter. Nilai \texttt{4} saat pemanggilan adalah argumen.
\end{pidkbox}
\end{frame}

\begin{frame}[fragile]{Parameter dan argumen}
\begin{lstlisting}
def luas_persegi_panjang(panjang, lebar):
    return panjang * lebar

A = luas_persegi_panjang(5, 3)
B = luas_persegi_panjang(lebar=3, panjang=5)
\end{lstlisting}
\begin{columns}[T,onlytextwidth]
\begin{column}{.48\linewidth}
\begin{pidkbox}{Positional argument}
Urutan argumen penting. Nilai pertama masuk ke parameter pertama.
\end{pidkbox}
\end{column}
\begin{column}{.48\linewidth}
\begin{pidkbox}{Keyword argument}
Nama parameter disebutkan langsung. Urutan tidak lagi menjadi fokus.
\end{pidkbox}
\end{column}
\end{columns}
\end{frame}

\begin{frame}[fragile]{Empat bentuk fungsi: bagian 1}
\begin{columns}[T,onlytextwidth]
\begin{column}{.48\linewidth}
\begin{pidkbox}{Tanpa argumen, tanpa \texttt{return}}
\begin{lstlisting}
def salam():
    print("Selamat belajar Python")
\end{lstlisting}
Cocok untuk menampilkan pesan atau instruksi.
\end{pidkbox}
\end{column}
\begin{column}{.48\linewidth}
\begin{pidkbox}{Tanpa argumen, dengan \texttt{return}}
\begin{lstlisting}
def konstanta_pi():
    return 3.14159
\end{lstlisting}
Cocok untuk menghasilkan nilai tetap.
\end{pidkbox}
\end{column}
\end{columns}
\end{frame}

\begin{frame}[fragile]{Empat bentuk fungsi: bagian 2}
\begin{columns}[T,onlytextwidth]
\begin{column}{.48\linewidth}
\begin{pidkbox}{Dengan argumen, tanpa \texttt{return}}
\begin{lstlisting}
def tampil_kuadrat(x):
    print(x**2)
\end{lstlisting}
Cocok untuk melaporkan hasil, bukan menghitung lanjut.
\end{pidkbox}
\end{column}
\begin{column}{.48\linewidth}
\begin{pidkbox}{Dengan argumen, dengan \texttt{return}}
\begin{lstlisting}
def kuadrat(x):
    return x**2
\end{lstlisting}
Cocok untuk menghasilkan nilai yang akan dipakai lagi.
\end{pidkbox}
\end{column}
\end{columns}
\end{frame}

\begin{frame}[fragile]{\texttt{print} tidak sama dengan \texttt{return}}
\begin{columns}[T,onlytextwidth]
\begin{column}{.48\linewidth}
\begin{lstlisting}
def g(x):
    return x**2 - 1

def h(x):
    print(g(x))
\end{lstlisting}
\end{column}
\begin{column}{.48\linewidth}
\begin{lstlisting}
y = h(3)
print(y)

z = h(3) + 2    # error
\end{lstlisting}
\end{column}
\end{columns}
\begin{pidkbox}{Ide utama}
\texttt{print} hanya menampilkan. Jika tidak ada \texttt{return}, fungsi mengembalikan \texttt{None}. Karena itu hasilnya tidak bisa langsung dipakai pada operasi aritmatika.
\end{pidkbox}
\end{frame}

\begin{frame}[fragile]{Fungsi dengan beberapa keluaran}
\begin{lstlisting}
def stat_ringkas(a, b, c, d):
    nilai = [a, b, c, d]
    minimum = min(nilai)
    maksimum = max(nilai)
    rata = sum(nilai) / len(nilai)
    rentang = maksimum - minimum
    return minimum, maksimum, rata, rentang

mn, mx, r, rg = stat_ringkas(4, 7, 1, 9)
\end{lstlisting}
\begin{pidkbox}{Apa yang terjadi?}
Python sebenarnya mengembalikan satu objek bertipe \texttt{tuple}. Unpacking membagi isi tuple ke beberapa variabel.
\end{pidkbox}
\end{frame}

\begin{frame}[fragile]{Tuple dan unpacking}
\begin{columns}[T,onlytextwidth]
\begin{column}{.48\linewidth}
\begin{lstlisting}
hasil = stat_ringkas(4, 7, 1, 9)
print(hasil)
print(type(hasil))
\end{lstlisting}
\end{column}
\begin{column}{.48\linewidth}
\begin{lstlisting}
minimum, maksimum, rata, rentang = hasil

# jumlah variabel harus cocok
x, y = hasil    # ValueError
\end{lstlisting}
\end{column}
\end{columns}
\begin{pidkbox}{Mengapa penting?}
Satu proses komputasi sering menghasilkan beberapa informasi: nilai fungsi, galat, jumlah iterasi, atau status konvergensi.
\end{pidkbox}
\end{frame}

\begin{frame}[fragile]{Scope: variabel lokal dan global}
\begin{lstlisting}
x = 10

def f(a):
    x = a + 2
    return x

print(f(5))
print(x)
\end{lstlisting}
\begin{columns}[T,onlytextwidth]
\begin{column}{.48\linewidth}
\begin{pidkbox}{Variabel lokal}
\texttt{x} di dalam fungsi hanya hidup di dalam fungsi tersebut.
\end{pidkbox}
\end{column}
\begin{column}{.48\linewidth}
\begin{pidkbox}{Variabel global}
\texttt{x} di luar fungsi tetap bernilai \texttt{10}. Nilai ini tidak otomatis berubah hanya karena ada \texttt{x} lokal.
\end{pidkbox}
\end{column}
\end{columns}
\end{frame}

\begin{frame}[fragile]{Mengapa variabel global perlu hati-hati?}
\begin{lstlisting}
total = 0

def tambah(x):
    global total
    total = total + x
    return total
\end{lstlisting}
\begin{pidkbox}{Masalah yang sering muncul}
\begin{itemize}
\item hasil fungsi bergantung pada riwayat pemanggilan;
\item pengujian menjadi tidak mandiri;
\item kesalahan sulit dilacak karena nilai dapat berubah dari tempat lain.
\end{itemize}
\end{pidkbox}
\begin{pidkbox}{Saran awal}
Utamakan fungsi yang menerima input secara eksplisit dan mengembalikan output secara eksplisit.
\end{pidkbox}
\end{frame}

\begin{frame}[fragile]{Fungsi sebagai alat dekomposisi}
\begin{columns}[T,onlytextwidth]
\begin{column}{.38\linewidth}
\[
 d(p,q)=\sqrt{(p_1-q_1)^2+(p_2-q_2)^2}
\]
\[
 E(p,q)=\frac{1}{1+d(p,q)^2}
\]
\end{column}
\begin{column}{.59\linewidth}
\begin{lstlisting}
def jarak(p, q):
    dx = p[0] - q[0]
    dy = p[1] - q[1]
    return (dx**2 + dy**2)**0.5

def energi(d):
    return 1 / (1 + d**2)

def energi_dua_titik(p, q):
    return energi(jarak(p, q))
\end{lstlisting}
\end{column}
\end{columns}
\end{frame}

\begin{frame}{Komposisi fungsi}
Jika \(g\) menerima \(x\), dan \(f\) menerima hasil dari \(g\), maka
\[
(f\circ g)(x)=f(g(x)).
\]
\begin{center}
\begin{tikzpicture}[node distance=1.0cm,>=Stealth]
\node[draw,rounded corners,fill=pidkblue,minimum width=2cm,minimum height=.85cm] (x) {$x$};
\node[draw,rounded corners,fill=pidknavy,text=white,minimum width=2.4cm,minimum height=.85cm,right=of x] (g) {$g$};
\node[draw,rounded corners,fill=pidkblue,minimum width=2.5cm,minimum height=.85cm,right=of g] (gx) {$g(x)$};
\node[draw,rounded corners,fill=pidknavy,text=white,minimum width=2.4cm,minimum height=.85cm,right=of gx] (f) {$f$};
\node[draw,rounded corners,fill=pidkblue,minimum width=2.8cm,minimum height=.85cm,right=of f] (fgx) {$f(g(x))$};
\draw[->,thick] (x)--(g); \draw[->,thick] (g)--(gx); \draw[->,thick] (gx)--(f); \draw[->,thick] (f)--(fgx);
\end{tikzpicture}
\end{center}
\begin{pidkbox}{Dalam Python}
Ide yang sama muncul saat satu fungsi memanggil fungsi lain. Program menjadi susunan unit kecil, bukan satu blok panjang.
\end{pidkbox}
\end{frame}

\begin{frame}[fragile]{Fungsi sebagai objek}
\begin{lstlisting}
def proses(x, f):
    return f(x)

def kuadrat(x):
    return x**2

def tambah1(x):
    return x + 1

print(proses(5, kuadrat))
print(proses(5, tambah1))
\end{lstlisting}
\begin{pidkbox}{Makna penting}
Nama \texttt{kuadrat} tanpa tanda kurung adalah objek fungsi. Objek itu dapat dikirim sebagai argumen ke fungsi lain.
\end{pidkbox}
\end{frame}

\begin{frame}[fragile]{List fungsi dan pipeline}
\begin{lstlisting}
def kuadrat(x):
    return x**2

def tambah1(x):
    return x + 1

def pipeline(x, ops):
    hasil = x
    for f in ops:
        hasil = f(hasil)
    return hasil

ops = [abs, kuadrat, tambah1]
print(pipeline(-3, ops))
\end{lstlisting}
\begin{pidkbox}{Baca alurnya}
Nilai \(-3\) diproses oleh \texttt{abs}, lalu \texttt{kuadrat}, lalu \texttt{tambah1}. Ini mirip alur transformasi data.
\end{pidkbox}
\end{frame}

\begin{frame}[fragile]{Lambda: fungsi pendek untuk operasi pendek}
\begin{columns}[T,onlytextwidth]
\begin{column}{.36\linewidth}
\[
 f(x)=3x^2-2x+1
\]
\begin{lstlisting}
f = lambda x: 3*x**2 - 2*x + 1
\end{lstlisting}
\end{column}
\begin{column}{.61\linewidth}
\begin{lstlisting}
xnilai = range(-3, 4)
y = [f(x) for x in xnilai]

z = list(map(lambda x: 3*x**2 - 2*x + 1,
             xnilai))
\end{lstlisting}
\end{column}
\end{columns}
\begin{pidkbox}{Gunakan secara wajar}
Lambda cocok untuk fungsi satu baris yang jelas. Jika logika mulai panjang, fungsi bernama lebih baik.
\end{pidkbox}
\end{frame}

\begin{frame}[fragile]{List comprehension untuk evaluasi fungsi}
\begin{lstlisting}
def f(x):
    return 3*x**2 - 2*x + 1

xnilai = [-3, -2, -1, 0, 1, 2, 3]
y = [f(x) for x in xnilai]
\end{lstlisting}
\begin{pidkbox}{Cara membaca}
Untuk setiap \texttt{x} di \texttt{xnilai}, hitung \texttt{f(x)}, lalu simpan hasilnya ke list baru.
\end{pidkbox}
\begin{pidkbox}{Catatan untuk M3}
Hari ini list dipakai sebagai wadah nilai secukupnya. Pembahasan list, dictionary, array NumPy, dan operasi matriks dibuat lebih formal pada pertemuan berikutnya.
\end{pidkbox}
\end{frame}

\begin{frame}{Rekursi: fungsi memanggil dirinya sendiri}
\begin{pidkbox}{Ide dasar}
Rekursi dipakai ketika masalah dapat dipecah menjadi masalah yang sama tetapi lebih kecil.
\end{pidkbox}
\begin{columns}[T,onlytextwidth]
\begin{column}{.48\linewidth}
\begin{pidkbox}{Base case}
Kasus paling sederhana yang jawabannya langsung diketahui. Tanpa base case, rekursi tidak berhenti.
\end{pidkbox}
\end{column}
\begin{column}{.48\linewidth}
\begin{pidkbox}{Recursive case}
Kasus umum yang memanggil fungsi yang sama untuk ukuran masalah yang lebih kecil.
\end{pidkbox}
\end{column}
\end{columns}
\vspace{2mm}
\[
\text{masalah ukuran } n \quad \longrightarrow \quad \text{masalah ukuran } n-1
\]
\end{frame}

\begin{frame}[fragile]{Contoh rekursi: faktorial}
\begin{columns}[T,onlytextwidth]
\begin{column}{.42\linewidth}
\[
 n!=n(n-1)!,\qquad 0!=1.
\]
\begin{align*}
5!&=5\cdot 4!\\
  &=5\cdot4\cdot3\cdot2\cdot1.
\end{align*}
\end{column}
\begin{column}{.55\linewidth}
\begin{lstlisting}
def faktorial(n):
    if n == 0:
        return 1
    return n * faktorial(n - 1)

print(faktorial(5))
\end{lstlisting}
\end{column}
\end{columns}
\begin{pidkbox}{Pertanyaan kelas}
Baris mana yang menjadi base case? Baris mana yang membuat masalah menjadi lebih kecil?
\end{pidkbox}
\end{frame}

\begin{frame}[fragile]{Rekursi untuk barisan}
Diberikan barisan
\[
 a_0=2,\qquad a_1=3,\qquad a_n=2a_{n-1}+a_{n-2} \quad (n\ge 2).
\]
\begin{lstlisting}
def a(n):
    if n == 0:
        return 2
    if n == 1:
        return 3
    return 2*a(n-1) + a(n-2)
\end{lstlisting}
\begin{pidkbox}{Latihan cepat}
Hitung manual \(a_2,a_3,a_4\). Setelah itu jalankan fungsi dan bandingkan.
\end{pidkbox}
\end{frame}

\begin{frame}[fragile]{Halt: kapan program berhenti?}
\begin{columns}[T,onlytextwidth]
\begin{column}{.48\linewidth}
\begin{pidkbox}{Berhenti}
\begin{lstlisting}
i = n
while i > 0:
    i = i // 2
\end{lstlisting}
Nilai \texttt{i} makin kecil dan akhirnya mencapai 0.
\end{pidkbox}
\end{column}
\begin{column}{.48\linewidth}
\begin{pidkbox}{Tidak berhenti untuk \(n>0\)}
\begin{lstlisting}
i = n
while i > 0:
    i = i + 1
\end{lstlisting}
Nilai \texttt{i} justru makin jauh dari kondisi berhenti.
\end{pidkbox}
\end{column}
\end{columns}
\begin{pidkbox}{Kaitannya dengan rekursi}
Rekursi juga harus bergerak menuju base case. Jika tidak, pemanggilan fungsi berlangsung terus sampai program gagal.
\end{pidkbox}
\end{frame}

\begin{frame}[fragile]{Kesalahan umum pada fungsi rekursif}
\begin{lstlisting}
def faktorial_salah(n):
    if n == 0:
        return 1
    return n * faktorial_salah(n)
\end{lstlisting}
\begin{pidkbox}{Diagnosis}
\begin{itemize}
\item base case ada, tetapi tidak pernah didekati;
\item pemanggilan ulang memakai \texttt{n}, bukan \texttt{n-1};
\item ukuran masalah tidak turun;
\item akibatnya fungsi tidak halt secara normal.
\end{itemize}
\end{pidkbox}
\end{frame}

\begin{frame}[fragile]{Merancang fungsi yang mudah diuji}
\begin{lstlisting}
def jarak_1d(x, y):
    return abs(x - y)

assert jarak_1d(3, 3) == 0
assert jarak_1d(3, 8) == 5
assert jarak_1d(-2, 3) == 5
\end{lstlisting}
\begin{pidkbox}{Kriteria awal}
\begin{itemize}
\item input jelas;
\item output jelas;
\item tidak bergantung pada variabel global;
\item dapat diuji dengan contoh yang jawabannya diketahui.
\end{itemize}
\end{pidkbox}
\end{frame}

\begin{frame}{Checklist membaca fungsi}
\begin{pidkbox}{Saat melihat fungsi baru, tanyakan lima hal ini}
\begin{enumerate}
\item Apa nama tugas yang dikerjakan fungsi?
\item Apa input yang dibutuhkan?
\item Apa output yang dikembalikan?
\item Apakah fungsi hanya mencetak, atau benar-benar mengembalikan nilai?
\item Apakah fungsi bergerak menuju kondisi berhenti jika memakai loop atau rekursi?
\end{enumerate}
\end{pidkbox}
\begin{pidkbox}{Untuk mahasiswa Matematika}
Biasakan menulis fungsi sebagai padanan komputasional dari fungsi, transformasi, norma, jarak, ringkasan, operator, atau barisan.
\end{pidkbox}
\end{frame}

\begin{frame}[fragile]{Latihan akhir 1: menulis fungsi}
\begin{pidkbox}{Kerjakan di kelas}
Buat fungsi berikut. Untuk setiap fungsi, tuliskan input, proses, dan output.
\end{pidkbox}
\begin{enumerate}
\item \texttt{nilai\_polinom(x, koef)} untuk menghitung \(a_0+a_1x+\cdots+a_nx^n\).
\item \texttt{stat\_ringkas(a,b,c,d)} yang mengembalikan minimum, maksimum, rata-rata, dan rentang.
\item \texttt{norm1(v)} untuk menghitung \(\|v\|_1=\sum_i |v_i|\) dari list angka.
\end{enumerate}
\end{frame}

\begin{frame}[fragile]{Latihan akhir 2: diagnosa kesalahan}
\begin{lstlisting}
def rata(data):
    total = 0
    for x in data:
        total = total + x
        return total / len(data)
\end{lstlisting}
\begin{pidkbox}{Pertanyaan}
\begin{enumerate}
\item Mengapa fungsi di atas hanya memproses elemen pertama?
\item Perbaiki posisi \texttt{return}.
\item Tambahkan dua \texttt{assert} untuk menguji fungsi.
\end{enumerate}
\end{pidkbox}
\end{frame}

\begin{frame}[fragile]{Latihan akhir 3: membuat fungsi rekursif}
Diberikan barisan
\[
 u_0=1,\qquad u_1=1,\qquad u_n=u_{n-1}+2u_{n-2}.
\]
\begin{enumerate}
\item Tulis fungsi rekursif \texttt{u(n)}.
\item Tentukan base case dan recursive case.
\item Hitung \texttt{u(5)} secara manual dan dengan Python.
\item Jelaskan mengapa fungsi tersebut halt untuk \(n\ge 0\).
\end{enumerate}
\end{frame}

\begin{frame}{Quiz e-learning M2}
\begin{pidkbox}{Materi quiz}
\begin{itemize}
\item membaca fungsi dan memprediksi output;
\item membedakan parameter dan argumen;
\item membedakan \texttt{print}, \texttt{return}, dan \texttt{None};
\item menelusuri scope lokal dan global;
\item menentukan hasil fungsi rekursif sederhana.
\end{itemize}
\end{pidkbox}
\begin{pidkbox}{Cara belajar}
Kerjakan prediksi dahulu di kertas atau catatan. Setelah itu jalankan Python untuk memeriksa dan memperbaiki pemahaman.
\end{pidkbox}
\end{frame}

\begin{frame}{Tugas mandiri M2}
\begin{pidkbox}{Yang dikumpulkan}
\begin{enumerate}
\item empat bentuk fungsi sesuai materi hari ini;
\item fungsi polinom dan fungsi statistik ringkas;
\item contoh scope lokal dan global beserta penjelasan;
\item fungsi \texttt{proses(x,f)} dan \texttt{pipeline(x,ops)};
\item satu fungsi rekursif untuk barisan yang Anda pilih.
\end{enumerate}
\end{pidkbox}
\begin{pidkbox}{Format}
Boleh menggunakan Google Colab, Jupyter, Spyder, VS Code, atau IDE Python lain. Kode harus dapat dijalankan ulang dan diberi komentar secukupnya.
\end{pidkbox}
\end{frame}

\begin{frame}{Penutup}
\begin{columns}[T,onlytextwidth]
\begin{column}{.48\linewidth}
\begin{pidkbox}{Hal yang perlu dikuasai}
\begin{itemize}
\item membuat fungsi dengan input dan output jelas;
\item memakai \texttt{return} secara tepat;
\item menelusuri scope;
\item memahami rekursi dan kondisi berhenti.
\end{itemize}
\end{pidkbox}
\end{column}
\begin{column}{.48\linewidth}
\begin{pidkbox}{Arah M3}
\begin{itemize}
\item representasi algoritma;
\item list, tuple, dictionary;
\item pseudocode;
\item array NumPy dan comprehension.
\end{itemize}
\end{pidkbox}
\end{column}
\end{columns}
\vspace{4mm}
\begin{center}
\large Fungsi yang baik membuat program lebih dekat dengan cara berpikir matematis: jelas inputnya, jelas transformasinya, jelas outputnya.
\end{center}
\end{frame}

\end{document}

